\subsection{Phenomenology, Astrophysics and Cosmology of Theories with Sub-Millimeter Dimensions and TeV Scale Quantum Gravity --- arXiv:hep-ph/9807344}

\textbf{Authors:} Nima Arkani-Hamed, Savas Dimopoulos, Gia Dvali

\paragraph{主な主張}
標準模型を3-brane上に局在させ重力を $n\geq2$ 個の大きな余剰次元へ伝播させる TeV-scale gravity の実験・天体物理・宇宙論的制限を系統的に調べ、$n>2$ は当時の制限を十分回避できる一方、$n=2$ では SN1987A 冷却や diffuse photon background から6次元 Planck scale におよそ $30\,\mathrm{TeV}$ 以上が必要であることを示した \cite{ArkaniHamed:1998ADDpheno}。

\paragraph{新規性}
large extra dimensions に伴う KK graviton の生成を中心に、低エネルギー実験・超新星・宇宙論を横断して TeV-scale quantum gravity の可行性を初期に体系化した点が新しい。

\paragraph{位置付け}
ADD 構想の現象論を確立した基礎論文であり、6次元・二余剰次元模型に対する重力 sector の標準的な制限を与える一方、bulk neutrino や $T^2$ の complex-structure dependence 自体は扱っていない。
